\documentclass[11pt]{article}
\usepackage[margin=1in]{geometry}
\usepackage{amsmath,amssymb,amsthm}
\usepackage{graphicx}
\usepackage{placeins}
\usepackage{booktabs}
\usepackage[expansion=false]{microtype}
\usepackage[hidelinks]{hyperref}
\setlength{\emergencystretch}{3em}

\theoremstyle{definition}
\newtheorem{proposition}{Proposition}

\title{\bf The Fermat-Type Dodge Ladder:\\
Primes $n^{2^k}+1$ on the $6N$ Skeleton and the\\
Non-Monotonic Richness%
\thanks{Code and measured data: \url{https://github.com/Ruqing1963/6N-dodge-ladder}}}
\author{Ruqing Chen\\
\small GUT Geoservice Inc., Montreal, Quebec, Canada\\
\small \texttt{ruqing@hotmail.com}}
\date{June 8, 2026}

\begin{document}
\maketitle

\begin{abstract}
Part~XXXII read prime richness as the product of per-prime ``dodges''---primes forbidden to divide a
sequence enhance it. We push that idea into higher degree along the Fermat-type family
$n^{2^k}+1$: $n^2+1$, $n^4+1$, $n^8+1$. The arithmetic is governed by a single cyclotomic fact:
$x^{2^k}\equiv-1\pmod q$ is solvable exactly when $q\equiv1\pmod{2^{k+1}}$, with $2^k$ roots. So
$n^{2^k}+1$ can be divided \emph{only} by $q=2$ and primes $q\equiv1\pmod{2^{k+1}}$, a vanishing
fraction $2^{-k}$ of all primes, while at those primes it is hit hard ($\omega=2^k$). Two findings,
both measured by direct (Miller--Rabin / BPSW) primality testing to $n\le10^{6}$ ($n^4+1$ up to
$10^{24}$). (i) The mod-$6$ \emph{wing} structure is degree-invariant: every $n^{2^k}+1$ splits its
primes right:left $=1{:}2$ exactly as $n^2+1$ does, because $n^{2^k}\equiv n^2\pmod6$ for all $k\ge1$;
the exponent is invisible to the coarse lattice. (ii) The Bateman--Horn richness $C=\prod_q g(q)$ is
\emph{non-monotonic} in $k$: as the degree doubles, the dodge widens (more enhancing primes) but the
hits deepen, and the two effects balance to a peak at $n^4+1$. We measure $C=1.376,\ 2.676,\ 2.075$ for
$k=1,2,3$---rising then falling---with the local product reproducing each to $\lesssim1\%$, including a
degree-eight Bateman--Horn check rarely attempted. The infinitude of primes $n^4+1$ is, like Landau's
problem, open; we make no claim about it. As throughout, this is a measurement, not a theorem.
\end{abstract}

\section{Introduction}
The richness principle of Part~XXXII \cite{PXXXII} was simple: a polynomial is prime-rich in proportion
to how many small primes its discriminant forbids from dividing it. For $n^2+1$ (Part~XXXI
\cite{PXXXI}) the forbidding rule was the character of $-1$: only primes $q\equiv1\pmod4$ can divide,
and the rest enhance. Raising the degree along $n^4+1$, $n^8+1$, and generally $n^{2^k}+1$ sharpens the
rule into a cyclotomic one and lets us watch the dodge tighten.

The governing fact is classical. The congruence $x^{2^k}\equiv-1\pmod q$ has a solution iff the order
of $-1$, namely $2$, is a $2^k$-th power in $(\mathbb{Z}/q)^\times$, which happens iff
$2^{k+1}\mid q-1$; equivalently $q$ splits completely in $\mathbb{Q}(\zeta_{2^{k+1}})$. Then there are
exactly $2^k$ roots. So $n^{2^k}+1$ is divisible only by $q=2$ and by primes $q\equiv1\pmod{2^{k+1}}$,
which are a fraction $2^{-k}$ of all primes by Dirichlet. The sequence dodges an ever-larger share of
the primes as $k$ grows---but where it is hit, it is hit on $2^k$ residues at once.

The family has a famous diagonal. At $n=2$ it is the Fermat numbers $F_k=2^{2^k}+1$, the sequence
whose first five terms ($3,5,17,257,65537$) led Fermat to guess they were all prime---until $F_5$ was
found divisible by $641$. The dodge ladder reframes that story: the early $F_k$ sit in a corridor whose
escaped-prime share grows like $1-2^{-k}$, so small terms meet very few possible divisors and look
prime; but the same exponent that widens the dodge sets the depth of any hit at $\omega=2^k$, and when
a divisor finally lands ($641\mid F_5$) it lands hard. Fermat saw the dodge; he did not see that the
penalty for being hit scales with it.

None of this is new as theory, and the infinitude of primes of the form $n^4+1$ is as open as Landau's
$n^2+1$. What we contribute is the $6N$ reading of the ladder, the observation that the wing split is
degree-invariant, and a direct high-precision measurement of the resulting richness---which turns out
to rise and then fall, peaking at the quartic.

\section{Setup: the cyclotomic dodge}
For $f(n)=n^{2^k}+1$ and a prime $q$ let $\omega(q)$ be the number of roots of $f\equiv0\pmod q$. By the
fact above,
\[
\omega(q)=\begin{cases}1,& q=2,\\ 2^k,& q\equiv1\pmod{2^{k+1}},\\ 0,& \text{otherwise.}\end{cases}
\]
The Bateman--Horn richness is $C=\prod_q g(q)$ with $g(q)=\bigl(1-\omega(q)/q\bigr)/(1-1/q)$, and
$Q(N):=\#\{n\le N:f(n)\ \text{prime}\}\sim C\int_2^N dt/\log(t^{2^k}+1)$. The local factor reads
\[
g(q)=\frac{q}{q-1}>1\quad(\text{$q$ dodged}),\qquad
g(q)=\frac{q-2^k}{q-1}<1\quad(q\equiv1\bmod2^{k+1},\ \text{$q$ hit}).
\]
Because the values $n^{2^k}+1$ grow like $n^{2^k}$, a root sieve would need primes to $N^{2^k}$, out of
reach; we therefore enumerate by direct primality testing---deterministic Miller--Rabin (twelve bases,
valid below $3.3\times10^{24}$, covering $n^4+1$ to $n\le10^{6}$) for $k\le2$, and the
Baillie--PSW test for $n^8+1$. The conditional product for $C$ converges by the Dirichlet balance of the
dodged ($g>1$) against the hit ($g<1$) primes, exactly as for $n^2+1$.

\section{The dodge tightens; the wings do not move}
Table~\ref{tab:ladder} and Figure~\ref{fig:fig1}(A) show the tightening. For $n^2+1$ the divisor primes
are $q\equiv1\pmod4$ ($5,13,17,29,\dots$); for $n^4+1$ only $q\equiv1\pmod8$ ($17,41,73,89,\dots$), so
$5,13,29,37$ are now dodged; for $n^8+1$ only $q\equiv1\pmod{16}$ ($17,97,113,\dots$), dodging
$41,73,89$ as well. Each step halves the share of primes that can divide.

The coarse geometry is untouched. Since squaring fixes the residues $\{0,1,3,4\}$ of $n^2\bmod6$, one
has $n^{2^k}\equiv n^2\pmod6$ for every $k\ge1$, so $n^{2^k}+1\bmod6$ is identical to $n^2+1\bmod6$:
odd $n$ are annihilated, $6\mid n$ sends the value to the right break-zone, $n\equiv2,4$ to the left.

\begin{proposition}
The wing split of $n^{2^k}+1$ is $1{:}2$ (right:left) for every $k\ge1$, independent of the degree.
\end{proposition}
\noindent We measure right:left $=0.500,\,0.499,\,0.517$ for $k=1,2,3$ (Fig.~\ref{fig:fig2}(B)); the
exponent is invisible to the $6N$ lattice. The character that the degree \emph{does} control lives
modulo $2^{k+1}$, not modulo $6$: Figure~\ref{fig:fig2}(A) shows the mod-$8$ spectrum of $n^4+1$, with
three-quarters of the primes enhancing and the $q\equiv1\pmod8$ quarter suppressed by $(q-4)/(q-1)$.

\begin{table}[t]
\centering\small
\begin{tabular}{ccccccc}
\toprule
$f$ & divides iff & $\omega$ & $N_{\max}$ & $Q$ & right:left & richness $C$ (meas.\ / prod.)\\
\midrule
$n^2+1$ & $q\equiv1\,(4)$  & $2$ & $10^6$ & $54{,}109$ & $0.500$ & $1.376\ /\ 1.373$\\
$n^4+1$ & $q\equiv1\,(8)$  & $4$ & $10^6$ & $52{,}610$ & $0.499$ & $2.676\ /\ 2.679$\\
$n^8+1$ & $q\equiv1\,(16)$ & $8$ & $10^5$ & $2{,}498$  & $0.517$ & $2.075\ /\ 2.093$\\
\bottomrule
\end{tabular}
\caption{The dodge ladder. Each doubling of degree restricts the divisor primes to a finer residue
class and deepens the hit ($\omega=2^k$). The wing split is degree-invariant; the richness $C$ rises
then falls. (The $n^8+1$ sample is smaller, so its measured $C$ carries a $\sim$2\% statistical error.)}
\label{tab:ladder}
\end{table}

\section{The non-monotonic richness}
The richness is a tug-of-war. Widening the dodge adds enhancing factors $q/(q-1)>1$ over a larger share
of primes, pushing $C$ up; deepening the hit replaces the modest $n^2+1$ suppression $(q-2)/(q-1)$ with
the severe $(q-2^k)/(q-1)$ on the surviving divisor primes, pulling $C$ down. The measured constants
(Table~\ref{tab:ladder}, Fig.~\ref{fig:fig1}(B)) show the balance tipping: $C$ climbs from $1.376$ at
$n^2+1$ to $2.676$ at $n^4+1$, then \emph{falls} to $2.075$ at $n^8+1$. The peak at the quartic is not
a product artifact: the direct count gives $C_{\rm emp}(n^4+1)=2.676$ against
$C_{\rm emp}(n^8+1)=2.075$, a $22\%$ drop far exceeding the $2\%$ sampling error. Throughout, the local
product reproduces the measured count to better than a percent---a Bateman--Horn verification carried,
for $n^8+1$, to degree eight, where values reach $10^{40}$ and direct primality testing is the only
route. The two ends of the ladder make the balance plain. As $k\to\infty$ the share of primes that can
divide, $2^{-k}$, tends to zero---almost every prime is dodged, an unbounded supply of enhancements---
yet at the vanishing minority of hit primes the value is split $2^k$ ways at once, a rarer and rarer but
ever deeper destruction. Neither extreme maximises $C$: the unbounded-dodge limit is throttled by the
unbounded-depth hits, and the two trends meet their best compromise at $k=2$, where the richness peaks.

\begin{figure}[t]
\centering
\includegraphics[width=\textwidth]{p33_fig1.pdf}
\caption{(A) The dodge map: the primes that can divide $n^{2^k}+1$ (coloured) thin out as the degree
doubles, restricted to $q\equiv1\pmod{2^{k+1}}$; dodged primes in grey. (B) The richness $C$ versus
degree: the product (blue) and the measured count (red) agree, and both peak at $n^4+1$ before falling
at $n^8+1$.}
\label{fig:fig1}
\end{figure}

\begin{figure}[t]
\centering
\includegraphics[width=\textwidth]{p33_fig2.pdf}
\caption{(A) The mod-$8$ quartic-character spectrum of $n^4+1$: three-quarters of primes are dodged
($g>1$), the $q\equiv1\pmod8$ quarter suppressed with $\omega=4$. (B) The $1{:}2$ wing split (right
$=1/3$) is identical for $n^2+1$, $n^4+1$, $n^8+1$: the $6N$ mask does not see the exponent.}
\label{fig:fig2}
\end{figure}

\FloatBarrier
\section{Scope and honest ledger}
The Fermat-type ladder shows one mechanism resolved into two moving parts. The exponent sets a
cyclotomic dodge---$n^{2^k}+1$ is divisible only by $q\equiv1\pmod{2^{k+1}}$, a fraction $2^{-k}$ of
primes, with $\omega=2^k$ where it is hit---and the resulting richness is non-monotonic, peaking at
$n^4+1$ because a wider dodge buys enhancement while a deeper hit costs suppression. The coarse $6N$
geometry, by contrast, is blind to the exponent: the wing split stays pinned at $1{:}2$.

We are explicit about the boundary. The solvability of $x^{2^k}\equiv-1\pmod q$ and its root count are
classical cyclotomic facts; the constant $C$ and the Bateman--Horn asymptotic are standard, and the
$n^2+1$ constant is Landau--Shanks. What this paper adds is the $6N$ reading of the ladder, the
degree-invariance of the wings, and a direct high-precision measurement of the richness---including the
non-monotonic peak and a degree-eight check. We make no claim about the infinitude of primes
$n^{2^k}+1$ for any $k\ge1$: that is open already for $k=1$ (Landau) and no finite computation can
settle it; the counts we report would look the same whatever the truth. Everything here equals a
direct primality count, takes the Hardy--Littlewood heuristic as input, and asserts no infinitude.

\paragraph{Data and code availability.} The Miller--Rabin / BPSW enumeration for $n^2+1$, $n^4+1$,
$n^8+1$, the local-factor products, the dodge and $\omega(q)$ tables, and the figure scripts are openly
available at \url{https://github.com/Ruqing1963/6N-dodge-ladder}.

\begin{thebibliography}{9}
\bibitem{HL} G.~H.~Hardy, J.~E.~Littlewood, \emph{Some problems of `Partitio Numerorum'; III}, Acta Math.\ \textbf{44} (1923), 1--70.
\bibitem{BH} P.~T.~Bateman, R.~A.~Horn, \emph{A heuristic asymptotic formula concerning the distribution of prime numbers}, Math.\ Comp.\ \textbf{16} (1962), 363--367.
\bibitem{Shanks} D.~Shanks, \emph{On the conjecture of Hardy \& Littlewood concerning the number of primes of the form $n^2+a$}, Math.\ Comp.\ \textbf{14} (1960), 321--332.
\bibitem{BPSW} R.~Baillie, S.~S.~Wagstaff, \emph{Lucas pseudoprimes}, Math.\ Comp.\ \textbf{35} (1980), 1391--1417.
\bibitem{PXXX} R.~Chen, \emph{Prime arithmetic progressions as a lattice resonance} (Part XXX), Zenodo, doi:10.5281/zenodo.20598405 (2026).
\bibitem{PXXXI} R.~Chen, \emph{Primes of the form $n^2+1$ on the 6N skeleton} (Part XXXI), Zenodo, doi:10.5281/zenodo.20600655 (2026).
\bibitem{PXXXII} R.~Chen, \emph{The richness spectrum of quadratic primes on the 6N skeleton} (Part XXXII), Zenodo, doi:10.5281/zenodo.20601304 (2026).
\bibitem{review} R.~Chen, \emph{Arithmetic geodynamics on the 6N skeleton: a review of Parts I--XIX}, Zenodo, doi:10.5281/zenodo.20585301 (2026).
\end{thebibliography}

\end{document}
